\section{Introduction}

The growing computational and memory demands of large language model (LLM) inference have driven the development of specialized hardware architectures. Among them, spatial dataflow architectures improve efficiency by streaming intermediate results directly between compute stages, reducing off-chip memory accesses~\cite{streamtensor}. FPGAs provide a flexible platform for implementing such workload-specific architectures. However, designing efficient FPGA spatial accelerators requires coordinated decisions on compute parallelism, stream buffering, and memory organization under resource and timing constraints~\cite{allo,streamtensor}. Moreover, evolving deployment requirements and physical implementation feedback often require revisiting earlier architectural decisions, leading to repeated engineering effort and long design cycles.

To reduce this design effort, compiler- and high-level synthesis (HLS)-based frameworks, such as Allo~\cite{allo}, HIDA~\cite{hida}, and StreamTensor~\cite{streamtensor}, generate spatial accelerators through structured intermediate representations (IRs) and compiler-defined optimization passes. These frameworks provide explicit representations and systematic mechanisms for hardware optimization. However, extending them to unsupported model structures, numerical formats, or hardware platforms may still require additional compiler and hardware development. Furthermore, their predefined transformation flows do not directly address the coordination of open-ended design revisions proposed by autonomous agents across heterogeneous hardware artifacts.

Meanwhile, LLM-based agents offer a more flexible approach by generating hardware code and iteratively refining designs using EDA tool feedback. Prior studies have explored RTL generation and debugging~\cite{mage,rtlfixer}, functional verification~\cite{assertllm,correctbench}, HLS optimization~\cite{hlspilot}, and multi-stage module design~\cite{mahl}. GPT4AIGChip~\cite{gpt4aigchip} and Design Conductor 2.0~\cite{designconductor2} further demonstrate accelerator-level automation. However, extending agentic automation to complete FPGA spatial accelerator development introduces a distinct difficulty: design decisions and revisions span model specifications, architectural configurations, generated RTL, and physical implementation, while their dependencies must remain consistent throughout the workflow.

The fundamental problem lies in maintaining a valid design state as interdependent hardware decisions evolve across stages. A design that passes individual stage-level checks may still violate constraints established elsewhere, while a change to one architectural decision may affect multiple components and previously obtained verification results. This nonlocal and revision-dependent nature of hardware design introduces two key challenges. \textbf{First, preserving cross-stage design consistency.} Early decisions, such as operator dimensions, numerical formats, and compute parallelism, constrain downstream hardware structures and interfaces. These constraints must be preserved during forward design and propagated to affected stages when earlier decisions are revised. Otherwise, locally valid components may become incompatible, particularly when LLM agents rely on incomplete or lengthy interaction histories~\cite{lostmiddle,longbench,locomo}. \textbf{Second, tracing downstream failures to their cross-stage root causes.} Failures observed during verification or physical implementation may originate from upstream architectural decisions rather than the reported component. For example, a timing violation may result from excessive compute parallelism instead of a local RTL defect. Without tracing the relevant dependencies, agents may repeatedly repair symptoms rather than underlying causes~\cite{swesearch}. Furthermore, such repairs may invalidate previous verification evidence, requiring affected components to be revalidated.

Addressing these challenges requires more than retaining design history: the workflow must track how decisions constrain one another, identify the components affected by design revisions, and establish which verification results remain valid. To this end, we propose \textbf{ACCORD}, a contract-guided multi-agent framework that organizes FPGA spatial accelerator development as a sequence of \emph{contract-preserving, evidence-backed design state transitions}. At iteration $t$, ACCORD maintains a shared design state $S_t=(G_t,A_t,R_t)$, where $G_t$ is a Design Contract Graph (DCG) capturing cross-stage design constraints and dependencies, $A_t$ contains the associated design artifacts, and $R_t$ records verification and implementation evidence bound to specific design revisions. Guided by this state, specialized agents construct spatial pipelines using reusable hardware templates and perform hardware-bound design space exploration. Each proposed revision is checked against affected contracts and tool requirements before being accepted. The same dependency information supports forward constraint propagation and backward failure diagnosis, while revision-bound evidence prevents outdated validation results from being reused. Together, these mechanisms enable consistent design refinement and evidence-backed design closure under the specified functional and implementation requirements.

The main contributions of this paper are summarized as follows:
\begin{itemize}
\item \textbf{Contract-preserving agentic design.}
We present ACCORD, an end-to-end FPGA spatial accelerator design framework that coordinates specialized agents through a shared design state and contract-preserving transitions, maintaining design decisions, artifacts, and validation evidence across stages.

\item \textbf{DCG-guided spatial accelerator generation.}
We introduce the DCG to capture cross-stage hardware constraints and dependencies, guiding template binding, spatial pipeline construction, hardware-bound design space exploration, and consistency maintenance during design revisions.

\item \textbf{Dependency-guided repair and design closure.}
We develop a failure diagnosis and repair mechanism that traces downstream failures to candidate upstream causes, restricts modifications to the affected scope, and maintains revision-bound verification evidence for reliable design closure.

\item \textbf{End-to-end evaluation.}
We evaluate ACCORD across LLM workloads, hardware configurations, and upstream LLM backbones, examining accelerator correctness and implementation quality, cross-stage consistency, failure diagnosis, and design efficiency.
\end{itemize}